% Time-stamp: <2025-02-26 18:15:25>


\documentclass[presentation,aspectratio=1610]{beamer}

%% Full theme: AnnArbor Antibes Bergen Berkeley Berlin Boadilla boxes CambridgeUS Copenhagen Darmstadt default Dresden EastLansing Frankfurt Goettingen Hannover Ilmenau JuanLesPins Luebeck Madrid Malmoe Marburg Montpellier PaloAlto Pittsburgh Rochester Singapore Szeged Warsaw
% \usetheme{Luebeck}

%% outer themes (header/footer): default infolines miniframes smoothbars sidebar split shadow tree smoothtree
% \useoutertheme[subsection=false]{smoothbars}
% \usetheme{Rochester}
 \useoutertheme{split}

%% inner theme (content): default circles rectangles rounded inmargin
\useinnertheme{rectangles}

\usecolortheme{beaver}

% Time-stamp: <2021-11-17 15:29:13 lperrin>


%!SECTION! Packages needed 

\usepackage{etex} % use when too many packages, vpv



\usepackage[english]{babel}

\usepackage{cmap}
\usepackage{graphicx}
\usepackage{color,colortbl,spot}
\usepackage{pgfplots}


% \usepackage{utopia}
% Possible fonts:
% mathptmx, helvet, bookman, chancery, charter, libertine, mathptm,
% newcent, palatino, pifont, utopia
% \usepackage[libertine]{newtxmath}
% \usepackage[scaled=0.96]{zi4}


\usepackage{hyperref}
\usepackage{appendixnumberbeamer}

\usepackage{booktabs} % tables, vpv
\usepackage{subcaption}
\usepackage{comment}


\usepackage{amsmath}
\usepackage{amssymb}


\usepackage{multicol}
\usepackage{multirow}
\usepackage{array} % To use monospace font in arrays easily
\newcolumntype{t}{>{\tt}c}



\usepackage[absolute,overlay]{textpos}
\usepackage{overpic}
\usepackage{tikz}
\usetikzlibrary{decorations.pathreplacing}
\usetikzlibrary{arrows}
\usetikzlibrary{arrows.meta}
\usetikzlibrary{patterns}
\usetikzlibrary{shapes}
\usetikzlibrary{shadows}
\usetikzlibrary{calc}
\usetikzlibrary{math}
\tikzstyle{fun}=[draw,very thick,fill=white]
\tikzstyle{fork}=[shape=circle,inner sep=0pt,minimum size=3pt,fill=black]
\tikzstyle{op}=[inner sep=2pt]
\tikzset{every picture/.style={line width=1.2pt}}



% !SECTION! Theme

% \usefonttheme{professionalfonts}

\usepackage{mathspec}
\setmainfont{Inria Serif}[Scale=MatchLowercase]
\setsansfont{Inria Sans}[Scale=MatchLowercase]
\setmathfont(Digits,Latin,Greek,Symbols)[Scale=MatchLowercase]{Inria Serif}
\setmathrm[Scale=MatchLowercase]{Inria Serif}


%!SUBSECTION! Custom colors 
\definecolor{light-gray}{RGB}{240, 240, 240}
\definecolor{dark-gray}{gray}{0.45}
\definecolor{navy}{RGB}{0, 40, 150}
\definecolor{darkgreen}{RGB}{10, 150, 60}
\definecolor{brown}{RGB}{120, 50, 10}
\definecolor{darkred}{rgb}{0.8,0,0}
\definecolor{mypink}{RGB}{200, 100, 100}



%% Inner color theme: lily orchid rose
%% Outer color theme: whale seahorse dolphin
%% Full  color theme: albatross beaver beetle crane dove fly seagull wolverine
% \usecolortheme{beaver}
% \usecolortheme{orchid}


%!SUBSECTION! Completely custom color theme 

\definecolor{mycol1}{RGB}{231, 51, 18}
\definecolor{mycol2}{RGB}{56,66,87}

\definecolor{mybrown}{RGB}{170,105, 57}
\definecolor{mydarkbrown}{RGB}{ 85, 36,  0}

\definecolor{red}{RGB}{192,0,0}
\definecolor{blue}{RGB}{ 50, 60,200}
\definecolor{teal}{RGB}{170,105, 57}

\setbeamercolor{section in toc}{fg=black,bg=white}
\setbeamercolor{alerted text}{fg=red}
\setbeamercolor*{palette primary}{fg=mycol2,bg=white}
\setbeamercolor*{palette secondary}{fg=mycol1,bg=mycol2}
\setbeamercolor*{palette tertiary}{bg=mycol2,fg=mycol1}
\setbeamercolor*{palette quaternary}{fg=white,bg=mycol2} % smoothbars

\setbeamercolor*{sidebar}{fg=darkred,bg=gray!15!white}

% \setbeamercolor*{palette sidebar primary}{fg=darkred!10!black}
% \setbeamercolor*{palette sidebar secondary}{fg=white}
% \setbeamercolor*{palette sidebar tertiary}{fg=darkred!50!black}
% \setbeamercolor*{palette sidebar quaternary}{fg=gray!10!white}


\setbeamercolor{titlelike}{fg=mycol1,bg=gray!15!white}
\setbeamercolor{frametitle}{fg=mycol1,bg=gray!15!white}
\setbeamercolor{frametitle right}{bg=mycol2!60!white}

\setbeamercolor{item}{fg=mycol1}
\setbeamercolor{subitem}{fg=mycol1}
\setbeamercolor{subsubitem}{fg=mycol1}

\setbeamercolor*{separation line}{}
\setbeamercolor*{fine separation line}{}

\setbeamercolor{block title}{fg=mycol1!80!black,bg=mycol1!10!white}
\setbeamercolor{block body}{bg=mycol1!2!white}
\setbeamercolor{block title example}{fg=mycol2!80!black,bg=mycol2!10!white}
\setbeamercolor{block body example}{bg=mycol2!5!white}

%!SUBSECTION! Setting footer 

\setbeamertemplate{footline}[frame number]
% \setbeamertemplate{footline}
% {
% \leavevmode%
% \hbox{%
% \begin{beamercolorbox}[wd=.3\paperwidth,ht=2.25ex,dp=1ex,center]{author in head/foot}%
%   \usebeamerfont{author in head/foot}\insertshortauthor
% \end{beamercolorbox}%
% \begin{beamercolorbox}[wd=.7\paperwidth,ht=2.25ex,dp=1ex,center]{title in head/foot}%
%   \usebeamerfont{title in head/foot}\insertshorttitle\hspace*{3em}
%   \hfill
%   \insertframenumber{} / \inserttotalframenumber\hspace*{1ex}
% \end{beamercolorbox}}%
% \vskip0pt%
% }
\beamertemplatenavigationsymbolsempty




%!SECTION! Abbreviations

\newcommand\smallcite[1]{{\color{gray} \footnotesize \cite{#1}}}


%!SECTION! Custom frame types

% Adding the plan at the beginning of each section starting from now
\newcommand\tocStartsAppearingHere{ %
  \begin{frame}[noframenumbering]
    \frametitle{Outline}
    \tableofcontents[
    sectionstyle=show,
    subsectionstyle=hide,
    subsubsectionstyle=hide] 
  \end{frame}

  \AtBeginSection[] { %
    % \begin{frame}[noframenumbering]
    %   \frametitle{Outline}
    %   \tableofcontents[
    %   currentsection,
    %   sectionstyle=show/shaded,
    %   subsectionstyle=hide,
    %   subsubsectionstyle=hide] 
    % \end{frame}

    \begin{frame}[noframenumbering]
      \frametitle{Plan of this Section}
      \tableofcontents[
      currentsection,
      currentsubsection,
      sectionstyle=show/shaded,
      subsectionstyle=show/hide/hide,
      subsubsectionstyle=hide]
    \end{frame}
  }

  \AtBeginSubsection[] { %
    \begin{frame}[noframenumbering]
      \frametitle{Plan of this Section}
      \tableofcontents[
      currentsection,
      currentsubsection,
      sectionstyle=show/shaded,
      subsectionstyle=show/shaded/hide,
      subsubsectionstyle=hide]
    \end{frame}
  }
}


\newcommand\myTitlePage{
  \frame[plain]{\titlepage \addtocounter{framenumber}{-1}}
}


%!SECTION! Useful macros
%=======================

% source: https://tex.stackexchange.com/questions/28704/defining-a-newcommand-with-variable-name-inside-another-newcommand
% usage: \defcolvar{f}{\mathcal{F}}{blue} creates new macro, \myf,
% which expands into {\color{blue}\mathcal{F}}
\newcommand{\defcolvar}[3]{%
  \expandafter\newcommand\csname my#1\endcsname{{\color{#3} #2 }} %
}

\newcommand\cols[4]{%
  \begin{columns}
    \begin{column}{#1\textwidth}
      \begin{center}
        #3
      \end{center}
    \end{column}
    \begin{column}{#2\textwidth}
      \begin{center}
        #4
      \end{center}
    \end{column}
  \end{columns}
}

\newcommand\showThenHighlight[1]{\only<1>{ #1 }
  \only<2>{\alert{\bfseries #1 }}}



%!SUBSECTION! Math
 
\newcommand\ftwo{\mathbb{F}_2}
\newcommand\proba{\mathsf{Pr}}

\newcommand\F{\mathbb{F}}
\newcommand\Zmod[1]{\mathbb{Z}/ #1 \mathbb{Z}}
\newcommand\field[1]{\mathbb{F}_{2^{#1}}}

\newcommand\openbutterfly[3]{\mathsf{H}^{#1}_{#2, #3}}
\newcommand\newbutterfly[1]{\mathsf{N}_{#1}}
\newcommand\linearspan[1]{\langle #1 \rangle}
\newcommand\rank{\mathsf{rank}}
\newcommand\extract[1]{\mathcal{X}_{#1}}
\newcommand\walshZeroes[1]{\mathcal{Z}_{#1}}
\newcommand\ddtZeroes[1]{\mathcal{Z}^D_{#1}}

\newcommand\lat[1]{\mathcal{W}_{#1}}
\newcommand\ddt[1]{\mathcal{D}_{#1}}
\newcommand\scalarprod[2]{#1 \cdot #2}

\newcommand\spaceInput{\mathcal{V}}
\newcommand\spaceOutput{\mathcal{V}^{\perp}}
\newcommand\allSpaces[1]{\mathcal{S}_{#1}}
\newcommand\admissible[1]{\mathcal{A}_{#1}}
\newcommand\invariant[1]{\mathcal{J}_{#1}}
\newcommand\codebook[1]{\Gamma_{#1}}
\newcommand\eaMappings{\mathcal{M}^{\textrm{EA}}}
\newcommand\nSpaces[1]{s_{#1}}
\newcommand\msb[1]{\mathsf{MSB}(#1)}

\newcommand\matSwap[1]{M_{#1}}
\newcommand\maj{\textrm{maj}}
\newcommand\twistable[1]{#1-twistable}
\newcommand\twotwoMat[4]{\left[ \begin{array}{cc} #1 & #2 \\ #3 & #4 \end{array} \right]}

\newcommand\tklog{TKlog}
\newcommand\tklogmath[1]{\mathcal{TK}_{#1}}

\newcommand\SBox{{\color{darkgreen}S}}
\newcommand\trace[1]{\textrm{Tr}\left( #1 \right)}
\newcommand\conference[1]{{\color{gray}(#1)}}

\newcommand\anomaly{\mathcal{A}}
\newcommand\symField[1]{\mathfrak{S}_{2^{#1}}}
\newcommand\tu[1]{\mathsf{TU}_{#1}}




\defcolvar{n}{n}{red}
\defcolvar{l}{\ell}{blue}

\title{Symmetric Techniques for Advanced Protocols: What *are* They?}
\author[Léo Perrin]{Léo Perrin\inst{1} }

\institute{Inria, Paris}
\titlegraphic{\includegraphics[height=1.1cm]{figures/inria}}


\date{14th of March 2025}




\begin{document}

{
  \pagestyle{empty}

  \maketitle


  \begin{frame}{Trendy topics}

    \begin{itemize}
      \large
    \item [] MPC-friendly? 
    \item [] Arithmetization-Oriented?
    \item [] Verification efficiency?
    \item [] Algebraic attacks? \pause
    \item [] Symmetric crypto \textbf{for the blockchain...} \pause
    \item [] \alert{... for neural networks???}
    \end{itemize}
    \begin{center}
      The conclusion of today: \textbf{symmetric cryptography} has always had to deal with specific \textbf{implementation criteria}, but the \alert{new ones} are indeed a bit \textbf{stranger than before}.
    \end{center}
  \end{frame}
}


\tocStartsAppearingHere{}


\section{What is the Purpose of a Symmetric Primitive}

\subsection{Let's look at primitives we all know}


\begin{frame}{Let's talk!}
  \vfill

  \begin{center}
    \includegraphics[width=8cm]{./figures/simpsons}
  \end{center}
  
  \vfill
\end{frame}


\begin{frame}{Unstable Definitions}
  \begin{alertblock}{What is ``efficient'' varies}
    \begin{itemize}
    \item What are the operations that we \textbf{can} use?
    \item What are the associated \textbf{costs}? 
    \end{itemize}
    \begin{center}
      How to get the best security for a given price?
    \end{center}
  \end{alertblock}

  
  
  \begin{exampleblock}{What is ``secure'' varies}
    \begin{itemize}
    \item Should the primitive work in many context? \hfill\onslide<4>{Modularity vs. Single use}
    \item Do we care about nonce-misuse? \hfill\onslide<4>{Robustness vs. ``not our problem''}
    \item
    \end{itemize}
    \begin{center}
      How do we define the security that the primitive must provide?
    \end{center}
  \end{exampleblock}
\end{frame}


\subsection{A Small Cog in a Big Machine}


\begin{frame}{Web Encryption}
  \begin{center}
    \begin{tikzpicture}[xscale=1.5,yscale=1.3]
      \draw[color=gray!50!white] (0, 0) rectangle (6, 4) node[pos=0.9,above] {Application};
      \onslide<2->{
        \draw[color=gray!70!white] (0, 0) rectangle (5.5, 3) node[pos=0.85,above] {Communications};
      }
      \onslide<3->{
        \draw[color=gray] (0, 0) rectangle (5, 2.5) node[pos=0.85,above] {\textbf{Secure} Library};
      }
      \onslide<4->{
        \draw[color=blue] (0, 0) rectangle (4.5, 2) node[pos=0.85,above] {Protocols};
      }
      \onslide<5->{
        \draw[color=red] (0, 0) rectangle (4, 1) node[pos=0.6,above] {Cryptographic Primitives};
      }
    \end{tikzpicture}

    \begin{itemize}
    \item<6-> We want \alert{software efficient} (computer and smartphone but not micro-controllers) efficient \alert{AEAD}.
    \item<7> AES-GCM; Chacha-poly1305.
    \end{itemize}
  \end{center}
\end{frame}



\begin{frame}{RAM Encryption}
  \begin{center}
    \begin{tikzpicture}[xscale=1.5,yscale=1.3]
      \draw[color=gray!50!white] (0, 0) rectangle (6, 4) node[pos=0.9,above] {Computer};
      \onslide<2->{
        \draw[color=gray!70!white] (0, 0) rectangle (5.5, 3) node[pos=0.85,above] {Components};
      }
      \onslide<3->{
        \draw[color=gray] (0, 0) rectangle (5, 2.5) node[pos=0.85,above] {\textbf{Secure} RAM};
      }
      \onslide<4->{
        \draw[color=blue] (0, 0) rectangle (4, 2) node[pos=0.85,above] {Trivial Protocol};
      }
      \onslide<5->{
        \draw[color=red] (0, 0) rectangle (4, 1) node[pos=0.6,above] {Cryptographic Primitives};
      }
    \end{tikzpicture}

    \begin{itemize}
    \item<6-> We want \alert{very low latency} \alert{block encryption}.
    \item<7> PRINCE? QARMA? not so clear at this stage.
    \end{itemize}
  \end{center}
\end{frame}


\begin{frame}{Some Constants}
  \begin{itemize}
    \setlength\itemsep{1cm}
    \large
  \item A symmetric primitive is very \alert{small} (but crucial) cog in a
    very big machine, \pause
  \item there are many \alert{different} ``big machines'', and \pause
  \item this has a \alert{huge influence} on what the primitive looks like.
  \end{itemize}
\end{frame}


\section{New Protocols}


\section{Conclusion}

\begin{frame}{Conclusion}
  \pause\vspace{1cm}
  
  \begin{center}
    \textbf{Thank you!}
  \end{center}
\end{frame}

\end{document}


% Local Variables:
% compile-command: "xelatex presentation.tex"
% End:
